\section{Back-end}

El desarrollo back-end corresponde a la capa del sistema encargada de gestionar la lógica de negocio, el procesamiento de datos y la comunicación con los sistemas gestores de bases de datos. Constituye el núcleo funcional de las aplicaciones web, ya que se encarga de atender las solicitudes del cliente, procesarlas y devolver las respuestas correspondientes de manera eficiente (Figura \ref{fig:Backend}).

Esta capa permite establecer una comunicación estructurada entre el cliente (generalmente el navegador) y el servidor, haciendo uso de lenguajes de programación y frameworks que facilitan la implementación de servicios, la gestión de recursos y la integración con otros sistemas. A través del back-end, se controlan aspectos fundamentales como la autenticación de usuarios, la validación de datos y la ejecución de operaciones sobre la información almacenada \parencite[]{BackendGG}.

\begin{figure}[H]
    \centering
    \includegraphics[width=0.5\linewidth]{teorico/backend.jpg}
    \caption{Flujo de un sistema web.}
    \label{fig:Backend}
\end{figure}

Asimismo, el desarrollo back-end desempeña un papel crucial en la seguridad e integridad de los sistemas, ya que de él depende la protección de los datos y la confiabilidad de las transacciones realizadas. Por ello, es fundamental considerar principios como la confidencialidad, la integridad y la disponibilidad de la información, así como la implementación de mecanismos de seguridad adecuados para resguardar los datos durante su almacenamiento y transmisión.

\subsection{.NET}

.NET es una plataforma de desarrollo de código abierto, gratuita y multiplataforma, desarrollada por Microsoft, que permite la creación de una amplia variedad de aplicaciones, incluyendo sistemas web, aplicaciones móviles, de escritorio, videojuegos y soluciones orientadas al Internet de las Cosas (IoT). Esta plataforma soporta múltiples lenguajes de programación, entre los que destacan C\# y Visual Basic, proporcionando un entorno robusto y versátil para el desarrollo de software (Figura \ref{fig:NET}).

\begin{figure}[ht]
    \centering
    \includegraphics[width=0.5\linewidth]{teorico/net.png}
    \caption{Áreas de utilidad de .NET}
    \label{fig:NET}
\end{figure}

El ecosistema de .NET se caracteriza por su amplia disponibilidad de herramientas que facilitan la integración de funcionalidades adicionales, promoviendo así un desarrollo más ágil, estructurado y mantenible. En el contexto del desarrollo web, .NET es ampliamente utilizado en la construcción del back-end, donde permite implementar la lógica de negocio, la gestión de datos y la exposición de servicios mediante arquitecturas modernas como APIs \parencite[]{Microsoft}.

Gracias a su enfoque en el rendimiento, la seguridad y la escalabilidad, .NET se posiciona como una solución sólida para el desarrollo del núcleo funcional de sistemas web, permitiendo la creación de aplicaciones eficientes y capaces de adaptarse a diferentes entornos y demandas.
